% ==================================================================
% id:       app-010
% topic:    שימושי הנגזרת
% subtopic: כלל לופיטל
% solution: full
% ==================================================================
חשבו את הגבולות הבאים:
\begin{solbox}
בחישוב הגבולות הבאים נשתמש בכלל לופיטל.
\end{solbox}
\bE\bMs{2}
\item $\ds{\lim_{x \to 1} \frac{x^x - x}{\ln x - x + 1}} $
\hint{זהו גבול מהצורה $\frac00$. כדי לגזור את $x^x$ כתבו אותו כ-$e^{x\ln x}$; ייתכן שתצטרכו להפעיל את כלל לופיטל יותר מפעם אחת.}
\begin{solbox}
לפני שניגש לגזירה, נבין מהי נגזרתה של
$h(x)=x^x$.
טוב, זה לא
$x$,
בחזקה קבועה וזה לא קבוע בחזקת
$x$
זהו "ערבוב" של שני אלה.
מה עשוי לעזור כאן?
אם רק היינו יכולים "להוריד" את המעריך, זה היה מאוד עוזר. מיהו ה"שפיץ" בהורדת המעריכים?
נכון! לוגריתם!
נתבונן בפונקציה
$$.h(x)=e^{\ln x^x}=e^{x\ln x}$$
בהסתמך על כלל השרשרת נקבל
$$,h'(x)=e^{x\ln x}(x\ln x)'=(\ln x+1)e^{x\ln x}=(\ln x+1)x^x.$$
הצלחנו!
ניגש לתרגיל המקורי.
נשים לב תחילה כי
$$,\lim_{x \to 1} \frac{x^x - x}{\ln x - x + 1}=\left[\frac{0}{0}\right]$$
מכיוון ש-
$f(x)=x^x - x,\ g(x)=\ln x - x + 1$
פונקציות גזירות בסביבה נקובה של
$x=1$
ו-
$g'(x)$
לא מתאפסת בסביבה זו,
אז נוכל לנסות להשתמש בכלל לופיטל.
ננסה לחשב את הגבול
$\lim_{x \to 1} \frac{f'(x)}{g'(x)}$
\begin{align*}
    \lim_{x \to 1} \frac{f'(x)}{g'(x)}=\ \lim_{x \to 1} \frac{\left(x^x - x\right)'}{(\ln x - x + 1)'}
    =\ \lim_{x \to 1}\frac{(\ln x+1)x^x-1}{\frac{1}{x}-1},
\end{align*}
לאחר הצבת
$x=1$
בביטוי האחרון נגלה ששוב קיבלנו גבול מסוג
$\frac{0}{0}$.
שוב, הפונקציות
$f'(x),\ g'(x)$
גזירות בסביבה נקובה של
$x=1$
ו-
$g''(x)\neq 0$
בסביבה זו,
נוכל לשתמש שוב בכלל לופיטל ולנסות לחשב את
$\displaystyle\lim_{x \to 1}\frac{f''(x)}{g''(x)}$.
מתקיים
\begin{align*}
    \lim_{x \to 1}\frac{f''(x)}{g''(x)}&=\ \lim_{x \to 1}\frac{\left((\ln x+1)x^x-1\right)'}{\left(\frac{1}{x}-1\right)'}\\
    &=\ \lim_{x \to 1}\frac{x^{x-1}+x^2(\ln x+1)}{-\frac{1}{x^2}}
    =\ -2.
\end{align*}
קיבלנו כי קיים הגבול
$\displaystyle\lim_{x \to 1} \frac{f''(x)}{g''(x)}$
לכן, לפי כלל לופיטל, קיים גם הגבול
$\lim_{x \to 1} \frac{f'(x)}{g'(x)}$
ושוב, לפי כלל לופיטל, קיים גם הגבול המקורי ומתקיים
$$.\lim_{x \to 1} \frac{\left(x^x - x\right)}{(\ln x - x + 1)}=-2$$
\end{solbox}

\item $\ds{\lim_{x \to +\infty} \frac{\ln x}{x^\eps} \quad (\eps > 0)}$
\begin{solbox}
נסמן
$f(x)=\ln x$,
$g(x)=x^\varepsilon$.
נשים לב כי
$f,g$
גזירות בכל הקרן החיובית ו-
$g'(x)=\varepsilon x^{\varepsilon-1}\neq 0$
לכל
$x>0$,
בנוסף
$$,\lim_{x \to +\infty} \frac{\ln x}{x^\varepsilon}=\left[\frac{\infty}{\infty}\right]$$
לכן ננסה להשתמש בכלל לופיטל,
כלומר נחשב את הגבול
$\displaystyle\lim_{x \to +\infty} \frac{f'(x)}{g'(x)}$
\begin{align*}
    \displaystyle\lim_{x \to +\infty} \frac{f'(x)}{g'(x)}=\ \lim_{x \to +\infty} \frac{(\ln x)'}{(x^\varepsilon)'}
    =\ \lim_{x \to +\infty}\frac{\frac{1}{x}}{\varepsilon \cdot x^{\varepsilon - 1}}
=\ \frac{1}{\varepsilon} \lim_{x \to +\infty} \frac{1}{x^\varepsilon}
=\ 0.
\end{align*}
מכיוון שקיים הגבול
$\displaystyle\lim_{x \to +\infty} \frac{f'(x)}{g'(x)}$,
נסיק מכלל לופיטל כי קיים גם הגבול המקורי ו-
$$.\displaystyle\lim_{x \to +\infty} \frac{f(x)}{g(x)}=0$$
\end{solbox}
\item $\ds{\lim_{x \to 0} \frac{\arcsin 2x - 2\arcsin x}{x^3}}$
\hint{הפעילו את כלל לופיטל, ולאחר מכן הוציאו מהגבול גורמים ששואפים לגבול סופי ושונה מאפס לפני שממשיכים.}
\begin{solbox}
נסמן
$.f(x)=\arcsin(2x)-2\arcsin x,\ g(x)=x^3$
שתי הפונקציות גזירות בסביבה נקובה של
$x=0$,
$g'(x)\neq0$
בסביבה זו ו-
$\displaystyle\lim_{x\to0}\frac{f(x)}{g(x)}=\left[\frac{0}{0}\right]$.
ננסה להשתמש בכלל לופיטל ונחשב את
$\displaystyle\lim_{x\to0}\frac{f'(x)}{g'(x)}$.
מתקיים
\begin{align*}
    \lim_{x \to 0} \frac{f'(x)}{g'(x)}&=\ \lim_{x \to 0} \frac{\frac{2}{\sqrt{1 - 4x^2}} - \frac{2}{\sqrt{1 - x^2}}}{3x^2}
=\ \frac{2}{3} \lim_{x \to 0} \frac{\sqrt{1 - x^2} - \sqrt{1 - 4x^2}}{x^2 \cdot \sqrt{1 - 4x^2} \cdot \sqrt{1 - x^2}}.
\end{align*}
נשים לב כי אם נוכל להתמודד עם הגבול
$\displaystyle\lim_{x\to0}\frac{\sqrt{1 - x^2} - \sqrt{1 - 4x^2}}{x^2}$,
נתמודד גם עם הגבול האחרון, הרי
$\displaystyle\lim_{x\to0}\sqrt{1 - 4x^2} \cdot \sqrt{1 - x^2}=1$.
נסמן
$$.h(x)=\sqrt{1 - x^2} - \sqrt{1 - 4x^2},\ k(x)=x^2$$
שוב, אנו נמצאים בתנאים של כלל לופיטל
(נמקו).
מתקיים
\begin{align*}
    \lim_{x\to0}\frac{h'(x)}{k'(x)}&=\ \lim_{x\to 0}\frac{-\frac{x}{\sqrt{1-x^2}}+\frac{4x}{\sqrt{1-4x^2}}}{2x}\\
    &=\ \lim_{x\to 0}\frac{-\frac{1}{\sqrt{1-x^2}}+\frac{4}{\sqrt{1-4x^2}}}{2}
    =\ \frac{-1+4}{2}
    =\ \frac{3}{2}.
\end{align*}
מכלל לופיטל, נסיק כי קיים הגבול
$\displaystyle\lim_{x\to0}\frac{h(x)}{k(x)}$
וערכו הוא
$\frac{3}{2}$.
מכאן הגבול המקורי הוא
$1$.
\end{solbox}
\item $\ds{\lim_{x \to 1} \left( 2 - x \right)^{\tan\frac{\pi x}{2}}}$
\hint{זהו גבול מהצורה $1^\infty$: כתבו את הביטוי כ-$e$ בחזקת מכפלה, והציגו את המכפלה כמנה שמתאימה לכלל לופיטל.}
\begin{solbox}
בדומה לתרגילים אחרים, ננסה להיפטר מהחזקה:
$$.\left( 2 - x \right)^{\tan\frac{\pi x}{2}}=e^{\ln\left( 2 - x \right)^{\tan\frac{\pi x}{2}}}=e^{\tan\frac{\pi x}{2}\ln\left( 2 - x \right)}=e^{\frac{\sin\frac{\pi x}{2}}{\cos\frac{\pi x}{2}}\ln\left( 2 - x \right)}$$
אם נדע להתמודד עם הגבול ב-
$x=1$
של
$\frac{1}{\cos\frac{\pi x}{2}}\ln\left( 2 - x \right)$,
נדע גם להתמודד עם הגבול
$\frac{\sin\frac{\pi x}{2}}{\cos\frac{\pi x}{2}}\ln\left( 2 - x \right)$
ב-
$x=1$,
הרי ש-
$$.\lim_{x\to1}\sin\frac{\pi x}{2}=1$$
נסמן
$$.f(x)=\ln\left( 2 - x \right),\ g(x)=\cos\frac{\pi x}{2}$$
מתקיים
$f,g$
גזירות בסביבה נקובה של
$x=1$,
$g'(x)\neq0$
בסביבה זו ו-
$$,\lim_{x\to0}\frac{f(x)}{g(x)}=\left[\frac{0}{0}\right]$$
לכן ננסה להשתמש בכלל לופיטל ונחשב את הגבול
$$.\lim_{x\to1}\frac{f'(x)}{g'(x)}$$
מתקיים
\begin{align*}
\lim_{x \to 1}\frac{f'(x)}{g'(x)}
&=\ \lim_{x\to 1}\frac{-\frac{1}{2-x}}{-\frac{\pi}{2}\sin\left(\frac{\pi x}{2}\right)}\\[4pt]
&=\ \frac{2}{\pi}.
\end{align*}
מכאן נסיק כי
$$.\lim_{x\to 1}(2-x)^{\tan\left(\frac{\pi x}{2}\right)}
=e^\frac{2}{\pi}$$
\end{solbox}

\item $\ds{\lim_{x \to 0} \lr{\cot x - \frac{1}{x}}}$
\hint{כתבו את $\cot x$ כמנה ואחדו את שני השברים למכנה משותף.}
\begin{solbox}
נציג תחילה
$$.\left(\cot x-\frac{1}{x}\right)
=\left(\frac{\cos x}{\sin x}-\frac{1}{x}\right)=\frac{x\cos x-\sin x}{x\sin x}=\frac{x\cos x-\sin x}{x^2\cdot\frac{\sin x}{x}}$$
אם נדע להתמודד עם הגבול ב-
$x=0$
של
$\frac{x\cos x-\sin x}{x^2}$,
,נדע גם להתמודד עם הגבול המקורי
הרי ש-
$$.\displaystyle\lim_{x\to0}\frac{\sin x}{x}=1$$
נסמן
$f(x)=x\cos x-\sin x,\ g(x)=x^2$.
מתקיים
$\displaystyle\lim_{x\to0}\frac{f}{g}=\left[\frac{0}{0}\right]$.
נשים לב כי שתי הפונקציות גזירות בסביבה נקובה של
$x=0$
ובסביבה זו
$g'(x)\neq0$.
לכן ננסה להשתמש בכלל לופיטל ולחשב את הגבול
$\displaystyle\lim_{x\to0}\frac{f'(x)}{g'(x)}$.
מתקיים
\begin{align*}
\displaystyle\lim_{x\to0}\frac{f'(x)}{g'(x)}
&=\ \lim_{x\to 0}\frac{\cos x-x\sin x-\cos x}{2x}\\
&=\ \lim_{x\to 0}\frac{-x\sin x}{2x}\\
&=\ \lim_{x\to 0}\frac{-\sin x}{2}\\\
&=\ 0.
\end{align*}
מכיוון שקיים הגבול
$\displaystyle\lim_{x\to0}\frac{f'(x)}{g'(x)}$,
קיים גם הגבול
$\displaystyle\lim_{x\to0}\frac{f(x)}{g(x)}$
וערכו הוא
$0$.
מכאן
\begin{align*}
    \lim_{x\to0}\left(\cot x-\frac{1}{x}\right)&=\ \lim_{x\to0}\frac{x\cos x-\sin x}{x^2\cdot\frac{\sin x}{x}}\\
    &=\ \lim_{x\to0}\frac{f(x)}{g(x)}\cdot\frac{x}{\sin x}\\
    &=\ 0\cdot1\\
    &=\ 0.
\end{align*}
\end{solbox}

\item $\ds{\lim_{x \to 0^+} \left( \ln \frac{1}{x} \right)^x}$
\hint{כתבו את הביטוי כ-$e^{x\ln\left(\ln\frac1x\right)}$, והציגו את המעריך (מהצורה $0\cdot\infty$) כמנה כדי להשתמש בכלל לופיטל.}
\begin{solbox}
נסדר תחילה את הפונקציה
$$.\left(\ln\frac{1}{x}\right)^x=e^{\ln\left(\ln\frac{1}{x}\right)^x}=e^{x\ln\left(\ln\frac{1}{x}\right)}$$
נשים לב כי
\begin{align*}
,\lim_{x\to 0^+}x\ln\left( \ln \frac{1}{x} \right)
&=\ [0\cdot\infty]
\end{align*}
נתאים את גבול לשימוש בכלל לופיטל ע"י הצגה
$$.x\ln\left( \ln \frac{1}{x} \right)=\frac{\ln\left( \ln \frac{1}{x} \right)}{\frac{1}{x}}$$
נסמן
$$.f(x)=\ln\left( \ln \frac{1}{x} \right),\ g(x)=\frac{1}{x}$$
נשים לב כי
$f,g$
גזירות בסביבה נקובה ימנית של
$x=0$
ו-
$g'(x)\neq0$
בסביבה זו וכן
$\displaystyle\lim_{x\to0^+}\frac{f(x)}{g(x)}=\left[\frac{\infty}{\infty}\right]$
, לכן ננסה להשתמש בכלל לופיטל.
נחשב
$\lim_{x\to 0^+}\frac{f'(x)}{g'(x)}$.

\begin{align*}
\lim_{x\to 0^+}\frac{f'(x)}{g'(x)}&= \lim_{x\to 0^+}\frac{\frac{1}{\ln(1/x)}\cdot\left(-\frac{1}{x}\right)}{-\frac{1}{x^2}}\\[6pt]
&=\ \lim_{x\to 0^+}\frac{x}{\ln(1/x)}\\
&=\ 0,
\end{align*}
מכיוון שקיים הגבול
$\displaystyle\lim_{x\to 0^+}\frac{f'(x)}{g'(x)}$,
קיים גם הגבול
$\displaystyle\lim_{x\to 0^+}\frac{f(x)}{g(x)}$
וערכו הוא
$0$.
מכאן
$$\lim_{x \to 0^+} \left( \ln \frac{1}{x} \right)^x=e^0=1$$
\end{solbox}
\eMs
\eE
